% !TEX TS-program = pdflatex
% !TEX encoding = UTF-8 Unicode

% This file is a template using the "beamer" package to create slides for a talk or presentation
% - Giving a talk on some subject.
% - The talk is between 15min and 45min long.
% - Style is ornate.

% MODIFIED by Jonathan Kew, 2008-07-06
% The header comments and encoding in this file were modified for inclusion with TeXworks.
% The content is otherwise unchanged from the original distributed with the beamer package.

\documentclass{beamer}
\setbeamercovered{invisible}



% Copyright 2004 by Till Tantau <tantau@users.sourceforge.net>.
%
% In principle, this file can be redistributed and/or modified under
% the terms of the GNU Public License, version 2.
%
% However, this file is supposed to be a template to be modified
% for your own needs. For this reason, if you use this file as a
% template and not specifically distribute it as part of a another
% package/program, I grant the extra permission to freely copy and
% modify this file as you see fit and even to delete this copyright
% notice. 


\mode<presentation>
{
  \usetheme{Warsaw}
  % or ...

  % or whatever (possibly just delete it)
}


\usepackage[english]{babel}
% or whatever

\usepackage[utf8]{inputenc}
\usepackage{qrcode}
\usepackage{hyperref}
% or whatever

\usepackage{times}
\usepackage[T1]{fontenc}
% Or whatever. Note that the encoding and the font should match. If T1
% does not look nice, try deleting the line with the fontenc.

%%% MATH RELATED

\title{MATH 180 Strategies of Problem Solving} 
\subtitle
{} % (optional)

\author[W.R. Casper] % (optional, use only with lots of authors)
{W.R. Casper}
% - Use the \inst{?} command only if the authors have different
%   affiliation.

\institute[California State University Fullerton] % (optional, but mostly needed)
{
  Department of Mathematics\\
  California State University Fullerton}
% - Use the \inst command only if there are several affiliations.
% - Keep it simple, no one is interested in your street address.

\subject{Talks}
% This is only inserted into the PDF information catalog. Can be left
% out. 



% If you have a file called "university-logo-filename.xxx", where xxx
% is a graphic format that can be processed by latex or pdflatex,
% resp., then you can add a logo as follows:

% \pgfdeclareimage[height=0.5cm]{university-logo}{university-logo-filename}
% \logo{\pgfuseimage{university-logo}}



% Delete this, if you do not want the table of contents to pop up at
% the beginning of each subsection:
\AtBeginSubsection[]
{
  \begin{frame}<beamer>{Outline}
    \tableofcontents[currentsection,currentsubsection]
  \end{frame}
}


% If you wish to uncover everything in a step-wise fashion, uncomment
% the following command: 

%\beamerdefaultoverlayspecification{<+->}


\begin{document}

\begin{frame}
  \titlepage
\end{frame}

\begin{frame}{Outline}
  \tableofcontents
  % You might wish to add the option [pausesections]
\end{frame}

% Since this a solution template for a generic talk, very little can
% be said about how it should be structured. However, the talk length
% of between 15min and 45min and the theme suggest that you stick to
% the following rules:  

% - Exactly two or three sections (other than the summary).
% - At *most* three subsections per section.
% - Talk about 30s to 2min per frame. So there should be between about
%   15 and 30 frames, all told.

\section{SoPS Lecture 2}
\subsection{Building conjectures}

\begin{frame}{Pinching Pennies}
  \begin{block}{Rules}
    \begin{itemize}
      \item arrange pennies in four piles
      \item turn: take one or more pennies \emph{but}
      \item take only from one pile per turn
      \item must take at least one penny
      \item person who takes last penny wins
    \end{itemize}
  \end{block}
\pause
  \begin{block}{Activity (20 mins)}
    \begin{itemize}
\pause
      \item split into pairs or groups of 3
\pause
      \item use small piles ($\leq 4$ coins per pile)
\pause
      \item {\color{red}record data:} initial board, who won? 
\pause
      \item try to find {\color{blue}patterns}
    \end{itemize}
  \end{block}
\end{frame}

\begin{frame}{Pinching Pennies: Debrief}
\pause
\textbf{Let's combine the data ...}
\pause
  \begin{block}{Question}
    What do the losing positions have in common?
  \end{block}
\pause
  \begin{block}{Submit your conjecture:}
    \begin{center}
      \qrcode[hyperlink,height=2cm]{https://padlet.com/williamrileycasper/pinching-pennies-losing-positions-h5xyswjj3ku3iqof}
    \end{center}
  \end{block}
\pause
  \begin{center}
    There seems to be some mysterious math behind the scenes...
  \end{center}
\end{frame}

\subsection{Expert play}

\begin{frame}{Vocabulary}
\pause
  \begin{block}{Binary number}
    A \textbf{binary number} is an expression like
      $$101_2,\quad001101_2,\quad 110100_2,\quad 1010101_2$$
  \end{block}
\pause
  \begin{block}{Binary vs decimal}
    \textbf{Binary numbers} are another way of writing natural numbers
      $$(b_r\dots b_3b_2b_1b_0)_2 = b_0 + 2b_1 + 4b_2 + 8b_3 + \dots + 2^rb_r$$
  \end{block}
\pause
  \begin{block}{Examples}
\pause
    $2 = 10_2 = 010_2 = 0010_2$\\
\pause
    $7 = 111_2,\quad 4 = 100_2, \quad 13 = 1101_2$
  \end{block}
\end{frame}

\begin{frame}{Try it yourself}
\pause
  \begin{block}{Question 1}
    What natural number is
      $$10110_2$$
  \end{block}
\pause
  \begin{block}{Question 2}
    What natural number is
      $$101101101_2$$
  \end{block}
\pause
  \begin{block}{Question 3}
    What binary number is the same as the natural number $2025$?
  \end{block}
\end{frame}

\begin{frame}{Observations}
\pause
  \begin{block}{Losing states}
    Some losing states for {\color{blue}Player 1}
    $$(3,2,1),\quad(4,6,2),\quad (3,3,0)$$
  \end{block}
\pause
  \begin{block}{Converted to binary}
    $$(011_2,010_2,001_2),\quad(100_2,110_2,010_2),\quad (011_2,011_2,000_2)$$
  \end{block}
\pause
  \begin{block}{Question}
    Any patterns?
  \end{block}
\end{frame}

\begin{frame}{Vocabulary}
\pause
  \begin{block}{binary XOR}
\pause
    The operation $\oplus$ called \textbf{binary XOR} adds binary numbers entrywise modulo $2$:
    $$110011_2 \oplus 101010_2 = 011001_2$$
  \end{block}
\pause
  \begin{block}{Losing states}
\pause
    $$(011_2,010_2,001_2)\rightarrow 011_2\oplus 010_2\oplus 001_2 = 000_2$$
\pause
    $$(100_2,110_2,010_2)\rightarrow 100_2\oplus 110_2\oplus 010_2 = 000_2$$
\pause
    $$(011_2,011_2,000_2)\rightarrow 011_2\oplus 011_2\oplus 000_2 = 000_2$$
  \end{block}
\end{frame}

\begin{frame}{Perfect play}
\pause
  \begin{block}{Theorem}
    If the binary XOR sum of the three piles of pennies is $0$, then it is a losing position for {\color{blue}Player 1}
  \end{block}
\pause
  \begin{block}{Proof}
    If the binary XOR sum is zero, the next move will make it not be zero.\\
\pause
    If the binary XOR sum is not zero, then there is a move to make it zero.\\
\pause
    Whatever {\color{blue}Player 1} does, he will leave the sum nonzero.\\
\pause
    Then {\color{red}Player 2} should make the sum zero again.\\
\pause
    Eventually we run out of stones, and that must happen on a turn when the sum becomes zero.
  \end{block}
\end{frame}

\begin{frame}{Wreck the professor}
  \begin{block}{Wreck the professor}
    Start with the initial piles $(5,7,9)$.
  \end{block}
  \begin{center}
    Fight me!
  \end{center}
\end{frame}

\begin{frame}{Break time!}
  \begin{block}{Activity (5 minutes)}
    \begin{itemize}
      \item get up, stretch get a cup of coffee
      \item chat up your neighbor, introduce yourself
    \end{itemize}
  \end{block}
\end{frame}

\begin{frame}{Practice makes perfect}
\pause
  \begin{block}{Activity (20 mins)}
    \begin{itemize}
\pause
      \item split into pairs or groups of 3
\pause
      \item use medium-sized piles ($7$-ish coins per pile)
\pause
      \item look at the starting state and figure out who will win with perfect play
\pause
      \item play the game through and see if you were right
    \end{itemize}
  \end{block}
\end{frame}

\subsection{Wrapping up}

\begin{frame}{Discussion}
\pause
  \begin{block}{Invariants}
    The binary XOR sum is an example of an \textbf{invariant}, something that one of the players can keep fixed. 
  \end{block}
\pause
  \begin{block}{Question 1}
    What did we learn today about invariants?
  \end{block}
\pause
  \begin{block}{Question 2}
    Faced with a new game, what strategy might we try using to determine perfect play?
  \end{block}
\end{frame}


\begin{frame}{Question of the Day}
\pause
  \begin{block}{Puzzle}
    You have $n$ pennies. On each turn, a player removes 1, 2, or 3 coins.  The player who takes the last coin wins. 
    For which $n$ is the start a winning position for {\color{blue}Player 1}?
  \end{block}
\pause
  \begin{block}{Activity (10 mins)}
    Discuss the puzzle with your neighbors.  What is the solution?
  \end{block}
\end{frame}

\begin{frame}{Final thoughts}
  \begin{block}{Reflection}
    \begin{itemize}
      \item What similarities and differences did you see between the "proof" we did today, vs. last class?
      \item Explain which player has a winning strategy in the previous puzzle and carefully describe what their strategy should be.
    \end{itemize}
  \end{block}
\begin{center}
Submit your responses by \textbf{Tonight!}
\end{center}
\end{frame}


\end{document}


